\hypertarget{native-533-lower}{}\label{native-533-lower}
\fancyhead[L]{\small KNITLMAYER ON PLANETARY DISTANCES.}
\begin{center}\small Göttingische gelehrte Anzeigen. Stück 179. pp. 1777..1782. November 10, 1810.\\[3pt]\rule{.72\linewidth}{.3pt}\end{center}

\emph{An Attempt at an Exact Account of the Progression in the Distances of Planets and Satellites from Their Central Bodies.} By \textsc{Ferdinand Knitlmayer}, captain in the Imperial and Royal Austrian Army. 1808. 48 pages in octavo. Brünn.

Several truths and relations in astronomy have, so to speak, been guessed from experience or by trial. They were later fully confirmed and led to important and far-reaching conclusions. These successes have encouraged many people to try their luck at comparing various numerical relations in our solar system and looking for a connection between them. This applies chiefly to the mean distances of the planets from the Sun. Insofar as the planets are regarded as individual bodies already formed, both their distances from the central body and all the other quantities that determine their orbits must indeed be regarded as entirely independent and arbitrary. There is no reason whatever to suppose that relations exist between them. If, however, we assume that our solar system owes its original formation to mechanical causes, we can certainly conceive of possible relations between the elements of the individual planets and, perhaps, their masses and densities. If such relations really could be discerned clearly, we might have some success in inferring from them how our planetary system first came into being. On this subject we have, in fact, only vague and inadmissible hypotheses. So far, all efforts of this kind, with which some lovers of astronomy have tormented themselves, have come to nothing. Hardly any mathematical astronomer will see in them more than an excusable but fruitless pastime.

For a time, people in Germany attached great weight to a supposed law into which they wanted to force the planetary distances. They treated them as sums of a constant and the terms of a geometric progression that doubles from one planet to the next. This gave only a rough approximation, without real agreement. Mercury had to be left out of the sequence altogether, since its distance was close to the constant itself and not to that sum. And no fewer than four claimants have been found for the place between Mars and Jupiter. Vesta, the last of them to become known, is by far the brightest and probably the largest of them all. Yet it is precisely Vesta that departs most markedly from the progression.

The purpose of this short work is to set out another law that the planetary distances are supposed to follow exactly rather than approximately. Mr. \textsc{Knitlmayer} starts with a geometric progression. He assumes that the actual mean distances of the individual planets depart farther from the terms of this progression as their masses increase and as the difference between their density and that of the Sun grows. The mean distances of the various planets are therefore to be represented by the formula
\SourceLeaf{534}
\[ a \,.\, 2^n + b\,m + c\,(D - \mathrm{1}) \]
where $m$ is the mass of each planet, $D$ its density with that of the Sun taken as the unit, and $a$, $b$, $c$ are constant coefficients. For $n$ the values are 1 for Mercury, 2 for Venus, 3 for the Earth, 4 for Mars, 5 for the new planets, 6 for Jupiter, 7 for Saturn and 8 for Uranus.

To test such a hypothesis, which in itself has no a priori basis and is, to call it by its proper name, pulled out of thin air, one would think that the values of the constants $a$, $b$, $c$ ought to be determined by elimination from the mean distances of the three planets whose masses and densities are best known, namely the Earth, Jupiter and Saturn. The formula could then be tested on planets whose masses and densities are known fairly accurately, namely Venus and Uranus. Since we really know the densities of Venus and Uranus only roughly, one would still have to be careful not to attach too much weight to any agreement that might turn up. The author, however, is far from proceeding in this way. He takes $a = \mathrm{0{,}075}$ straight away, for no apparent reason other than that the same geocentric progression had been used as the basis of the hypothesis mentioned above. He determines the coefficients $b$ and $c$ from the mean distances of the Earth and Jupiter. If we recast his presentation, which is for the most part very confused, in our own form, the results are $b = \mathrm{427{,}435}$, $c = \mathrm{0{,}135143}$ (taking the mass of the Sun $= \mathrm{1}$). For Venus, Mars, Saturn and Uranus he then takes the masses given by \textsc{Wurm} and \textsc{Laplace} and uses his hypothetical formula to calculate the densities from them. From these and the masses he calculates the volume and the apparent semidiameter at distance 1. He takes the fairly close agreement of his calculated values with the observed ones as proof that his hypothesis is correct. He then has little difficulty deducing the otherwise unknown masses and densities of Mercury and the four new planets from the apparent diameters determined by \textsc{Schröter}. At first sight, this supposed agreement might dazzle some readers. We shall therefore simply describe how the author goes about his calculations. For Vesta, Saturn and Uranus the above formula gives $D - \mathrm{1}$ a negative value $= -p$, so the author declares these planets less dense than the Sun. Yet instead of setting the density $= \mathrm{1} - p$, he calculates as if it were $= -\dfrac{1}{1+p}$. After this sample of the author’s mathematical knowledge, we will probably be excused from making any further remarks.

The author now applies his hypothesis to the satellite systems of the primary planets as well. We describe the method as follows, drawing out his real idea from the actual applications he gives, since his own account is in part incomprehensible and contradictory. He assumes that the mean distances of the individual satellites belonging to a single primary planet must likewise be represented by the formula
\[ \mathrm{0{,}075} \,.\, 2^n + \mathrm{427{,}435}\, m + \mathrm{0{,}135143}\,(D - \mathrm{1}) \]
(where $m$ is the mass, $D$ the density and $n$ the ordinal number of each satellite). The mean distance of the Earth from the Sun, used as the unit of length above, is to be replaced by a different unit for each planet. The author’s rule sets this unit at $\dfrac{A\BareRoot[3]{M}}{B}$, where $A$ denotes the mean distance of Mercury from the Sun, $B$ the mean distance from the Sun of the primary planet that the satellite accompanies, and $M$ the mass of that planet (always taking the mass of the Sun $= \mathrm{1}$). The author seems to have chosen this value for the unit of length so as to obtain \SourceLeaf{535}fair agreement for the Earth’s satellite. This value, however, has no other theoretical basis, and the satellites of the other planets cannot be used to test the hypothesis since we have no knowledge at all of their densities. We can therefore see in this whole application of the hypothesis to satellite systems nothing more than an empty game with numbers.
